% Material para Aristotle: secoes 2, 3, Definicao 4.1, Teorema 4.2 e secoes 5 e 6. Sem arquivos auxiliares.
% Enunciados e provas matematicos transcritos da fonte, sem alteracao. Figuras externas apenas por legendas.
% Figuras apresentadas por legendas apenas. Bibliografia externa omitida.
\documentclass[11pt,a4paper]{article}
\usepackage[T1]{fontenc}
\usepackage[utf8]{inputenc}
\usepackage{amsmath,amssymb,amsthm,mathtools}
\usepackage{booktabs,enumitem}
\usepackage[round,authoryear]{natbib}
\newtheoremstyle{preciso}{8pt}{5pt}{\itshape}{}{\bfseries}{.}{.5em}{}
\theoremstyle{preciso}
\newtheorem{teorema}{Teorema}[section]
\newtheorem{lema}[teorema]{Lema}
\newtheorem{proposicao}[teorema]{Proposição}
\newtheorem{corolario}[teorema]{Corolário}
\theoremstyle{definition}
\newtheorem{definicao}[teorema]{Definição}
\newtheorem{algoritmo}[teorema]{Algoritmo}
\newtheorem{observacao}[teorema]{Observação}
\renewcommand{\proofname}{Demonstração}
\newcommand{\N}{\mathbb N}
\newcommand{\Z}{\mathbb Z}
\newcommand{\Q}{\mathbb Q}
\newcommand{\R}{\mathbb R}
\newcommand{\D}{\mathbb D}
\newcommand{\LL}{\mathcal L}
\newcommand{\Asm}{\mathcal A}
\newcommand{\FF}{\mathcal F}
\newcommand{\EE}{\mathcal E}
\newcommand{\RR}{\mathcal R}
\newcommand{\CC}{\mathcal C}
\newcommand{\heightq}{\operatorname{ht}}
\newcommand{\rank}{\operatorname{rank}}
\newcommand{\Eval}{\operatorname{Eval}}
\newcommand{\ser}{\operatorname{ser}}
\newcommand{\norm}[1]{\left\|#1\right\|_{\infty}}
\newcommand{\normK}[2]{\left\|#1\right\|_{\infty,#2}}
\newcommand{\eqn}[2]{\begin{equation}\label{eq:#1}#2\end{equation}}
\newcommand{\figura}[4]{\par\refstepcounter{figure}\label{fig:#4}\noindent\textit{Figura \thefigure~(ilustracao omitida), legenda original: #3}\par}
\title{Biblioteca determinística de segmentos afins racionais em células diádicas\\
Taxas quantitativas e refinamento adaptativo — Fase 3 para formalização}
\author{Victor Gabriel de Oliveira}
\date{}
% OBJETIVO PARA ARISTOTLE: Secoes 5 e 6; seções 2, 3 e inicio da secao 4 fornecem dependencias.
% Preservar os enunciados; nao certificar algoritmos cuja semantica nao foi formalizada.
\begin{document}
\maketitle
\setcounter{section}{1}

% --- INICIO SECAO 2 (transcricao literal) ---
\section{Ferramentas de indexação}
Escreva $\N_0=\{0,1,2,\ldots\}$ e $T_m=m(m+1)/2$.

\subsection{Pareamento de Cantor}
Para $u,v\in\N_0$, defina
\eqn{01}{\pi(u,v)=T_{u+v}+v=\frac{(u+v)(u+v+1)}2+v.}
\begin{lema}[pareamento e cota]
$\pi:\N_0^2\to\N_0$ é bijetiva. Para $n\in\N_0$, sua inversa é
\eqn{02}{w=\left\lfloor\frac{\sqrt{8n+1}-1}{2}\right\rfloor,\qquad v=n-T_w,\qquad u=w-v.}
Além disso, $\pi(u,v)\le(u+v+1)^2$.
\end{lema}
\begin{proof}
Se $m=u+v$, então $T_m\le\pi(u,v)=T_m+v<T_{m+1}$, pois $0\le v\le m$. Assim, cada imagem pertence a um único bloco diagonal $[T_m,T_{m+1})$. Dado $n$, a expressão para $w$ em \eqref{eq:02} é precisamente o único inteiro com $T_w\le n<T_{w+1}$. Logo $0\le v=n-T_w\le w$, $u=w-v\ge0$ e $\pi(u,v)=n$, provando sobrejetividade. O bloco determina $w=u+v$ e, dentro dele, $v=n-T_w$; portanto o par é único. Por fim, com $x=u+v+1$, tem-se $T_{u+v}\le x^2/2$ e $v\le x^2/2$, donde $\pi(u,v)\le x^2$.
\end{proof}

\subsection{Inteiros e racionais por altura}
Defina $Z:\N_0\to\Z$ por
\eqn{03}{Z(t)=\begin{cases}t/2,&t\text{ par},\\-(t+1)/2,&t\text{ ímpar}.\end{cases}}
\begin{lema}
$Z$ é bijetiva, com $Z^{-1}(m)=2m$ para $m\ge0$ e $Z^{-1}(m)=-2m-1$ para $m<0$. Em particular, $Z^{-1}(m)\le2|m|+1$.
\end{lema}
\begin{proof}
Os índices pares percorrem $\N_0$ exatamente uma vez; os ímpares percorrem $\Z\setminus\N_0$ exatamente uma vez. As duas imagens são disjuntas e as fórmulas inversas resultam por resolução direta de \eqref{eq:03}.
\end{proof}

Todo $r\in\Q$ tem uma representação reduzida única $p/q$, com $p\in\Z$, $q\in\N$ e $\gcd(|p|,q)=1$. Adote a altura multiplicativa elementar
\eqn{04}{\heightq(p/q)=\max\{|p|,q\}.}
Ordene os pares reduzidos primeiro por altura, depois por $q$ e finalmente por $p$. Essa leitura define bijeções
\eqn{05}{Q:\N_0\to\Q,\qquad Q^*:\N_0\to\Q^*,}
na segunda das quais se omite $0/1$.

\begin{lema}[contagem por altura]
Seja $\Q_H=\{r\in\Q:\heightq(r)\le H\}$. Então
\eqn{06}{|\Q_H|\le3H^2\qquad\text{para todo }H\ge1,}
e existem $H_0$ e $c_1>0$ universais tais que
\eqn{07}{|\Q_H|\ge c_1H^2\qquad\text{para todo }H\ge H_0.}
Pode-se tomar $c_1=1/4$.
\end{lema}
\begin{proof}
A forma reduzida com denominador positivo é única: de $p/q=p'/q'$ e das duas coprimalidades segue $q\mid q'$ e $q'\mid q$, logo $q=q'$ e $p=p'$. Cada $r\in\Q_H$ corresponde então a um único par $(p,q)$ com $|p|\le H$ e $1\le q\le H$; há no máximo $(2H+1)H\le3H^2$ tais pares, o que prova \eqref{eq:06}.

Para \eqref{eq:07}, observe que $\Q_H$ contém todas as frações $p/q$ com $1\le q\le H$, $1\le p\le q$ e $\gcd(p,q)=1$, cuja quantidade é $\Phi(H)=\sum_{q\le H}\varphi(q)$. Pelo teorema clássico de Mertens sobre a soma de Euler \citep[Teorema 330]{hardy2008},
\[
\Phi(H)=\frac3{\pi^2}H^2+O(H\log H),
\]
e $3/\pi^2>0{,}30$. Logo $\Phi(H)\ge H^2/4$ para todo $H$ suficientemente grande.
\end{proof}

\begin{lema}[posto por altura]
As aplicações em \eqref{eq:05} são bijetivas. Se $\heightq(p/q)\le H$, então o posto de $p/q$ em qualquer uma delas é estritamente menor que $3H^2$.
\end{lema}
\begin{proof}
A ordenação esgota cada casca finita de altura antes de passar à seguinte; portanto lista todo racional uma única vez, o que dá a bijetividade. Todo elemento que precede $p/q$ pertence a $\Q_H$, e o lema~2.3 limita essa cardinalidade por $3H^2$. A remoção de $0/1$ em $Q^*$ só diminui a contagem.
\end{proof}

\subsection{Bijetividade e crescimento do índice}
A ordem de Calkin--Wilf fornece uma comparação entre bijetividade e crescimento do endereço. Seja $(S_n)_{n\ge0}$ a sequência diatômica de Stern,
\eqn{08}{S_0=0,\qquad S_1=1,\qquad S_{2n}=S_n,\qquad S_{2n+1}=S_n+S_{n+1},}
e $CW(n)=S_n/S_{n+1}$, $n\ge1$, que enumera $\Q_{>0}$ sem repetições.
\begin{lema}
Para todo $q\ge1$, $CW(2^{q-1})=1/q$.
\end{lema}
\begin{proof}
Da recorrência par, $S_{2^m}=S_1=1$. Também $S_{2^m+1}=m+1$: o caso $m=0$ é $S_2=1$ e
\[
S_{2^{m+1}+1}=S_{2(2^m)+1}=S_{2^m}+S_{2^m+1}=1+(m+1).
\]
Tomando $m=q-1$, resulta $CW(2^{q-1})=1/q$.
\end{proof}

O racional $1/q$ tem altura $q$, mas índice $2^{q-1}$; a escrita binária do índice exige $q$ bits, contra os $O(\log q)$ do lema~2.4. A ordenação por altura controla o endereço em função de $\heightq(r)$. A seção~7 determina o comprimento mínimo no pior caso por contagem e constrói uma enumeração que o realiza.
\figura{01_calkin_wilf}{.94}{Primeiros níveis da árvore de Calkin--Wilf. O ramo esquerdo contém $1/q$ no índice $2^{q-1}$; a ordenação por altura fornece índices polinomiais na altura.}{cw}

% --- INICIO SECAO 3 (transcricao literal) ---
\section{Células diádicas e operador gerador}
\subsection{Células, posse e centros}
Para $k\in\N_0$ e $m\in\Z$, sejam
\eqn{09}{J_{k,m}=[m2^{-k},(m+1)2^{-k}],\qquad J_{k,m}^{\rm own}=[m2^{-k},(m+1)2^{-k}),}
e
\eqn{10}{x_0(k,m)=\left(m+\frac12\right)2^{-k}=\frac{2m+1}{2^{k+1}}.}
A célula fechada sustenta o segmento; a versão semiaberta atribui cada fronteira a um único proprietário. Escreva
\[
\D=\{j2^{-K}:j\in\Z,\ K\in\N_0\},\qquad \D_{\rm mid}=\D\setminus\Z.
\]
\begin{lema}[estrutura diádica]
Para cada $k$, as células $J_{k,m}^{\rm own}$ particionam $\R$ e
\eqn{11}{J_{k,m}=J_{k+1,2m}\cup J_{k+1,2m+1},\qquad |J_{k,m}|=2^{-k}.}
O mapa $(k,m)\mapsto x_0(k,m)$ é uma bijeção de $\N_0\times\Z$ sobre $\D_{\rm mid}$, que é denso em $\R$.
\end{lema}
\begin{proof}
Para $x\in\R$, o único índice de posse na escala $k$ é $m=\lfloor2^kx\rfloor$; isso prova a partição. A identidade de refinamento decorre de cindir $[2m2^{-(k+1)},(2m+2)2^{-(k+1)}]$ em seu ponto médio.

Pela \eqref{eq:10}, o centro tem numerador ímpar e denominador reduzido $2^{k+1}\ge2$, logo pertence a $\D_{\rm mid}$. Reciprocamente, cada $x\in\D_{\rm mid}$ possui forma reduzida única $j/2^K$, com $K\ge1$ e $j$ ímpar; tomando $k=K-1$ e $m=(j-1)/2$, obtém-se $x=x_0(k,m)$. A mesma unicidade prova a injetividade. Por fim, os centros de escala $k$ estão espaçados por $2^{-k}$; escolhendo o mais próximo de um real dado, a distância tende a zero com $k$, o que prova a densidade.
\end{proof}
\figura{02_celulas}{.92}{Células diádicas, centros e posse das fronteiras. Cada célula se divide em dois filhos; os centros determinam a célula de origem.}{celulas}

\subsection{Retas ancoradas e descompactação}
Dados $x_0,y_0\in\R$ e $s\ne0$, a reta ancorada é
\eqn{12}{\ell(x)=y_0+s(x-x_0)=a+bx,\qquad b=s,\quad a=y_0-sx_0.}
Para $n\in\N_0$, aplique sucessivamente a inversa do pareamento:
\eqn{13}{(u,v)=\pi^{-1}(n),\qquad(k,\mu)=\pi^{-1}(u),\qquad(r,t)=\pi^{-1}(v).}
\begin{lema}
O mapa $n\mapsto(k,\mu,r,t)$ de \eqref{eq:13} é uma bijeção de $\N_0$ sobre $\N_0^4$.
\end{lema}
\begin{proof}
Sua inversa explícita é
\[
(k,\mu,r,t)\longmapsto\pi\bigl(\pi(k,\mu),\pi(r,t)\bigr),
\]
uma composição de bijeções pelo lema~2.1.
\end{proof}
Defina
\eqn{14}{m=Z(\mu),\quad x_{0,n}=x_0(k,m),\quad y_{0,n}=Q(r),\quad s_n=Q^*(t),}
e converta para a forma global por
\eqn{15}{b_n=s_n,\qquad a_n=y_{0,n}-s_nx_{0,n}.}
Se $x_L=m2^{-k}$ e $x_R=(m+1)2^{-k}$, ponha
\eqn{16}{c_n=\min\{\ell_n(x_L),\ell_n(x_R)\},\qquad d_n=\max\{\ell_n(x_L),\ell_n(x_R)\}.}
\begin{definicao}[gerador e biblioteca]
Se $\mathcal I$ é o conjunto dos intervalos fechados não degenerados de $\R$, defina
\eqn{17}{\Gamma:\N_0\to\R\times(\R\setminus\{0\})\times\mathcal I,\qquad
\Gamma(n)=(a_n,b_n,[c_n,d_n]).}
A biblioteca é $\LL=\Gamma(\N_0)$. O átomo de endereço $n$ é a restrição de $\ell_n$ à célula $J_{k(n),m(n)}$.
\end{definicao}
Como $s_n\ne0$ e $x_L\ne x_R$, o intervalo em \eqref{eq:16} é não degenerado. A inclinação não nula torna a interface invertível no eixo espacial: uma peça constante teria imagem pontual e perderia a informação da célula. Para $\Gamma(n)=(a_n,b_n,[c_n,d_n])$, a pré-imagem espacial é
\eqn{18}{I_n^x=\left[
\min\left\{\frac{c_n-a_n}{b_n},\frac{d_n-a_n}{b_n}\right\},
\max\left\{\frac{c_n-a_n}{b_n},\frac{d_n-a_n}{b_n}\right\}\right].}

\subsection{Recuperação da célula e endereço canônico}
\begin{lema}[validade espacial exata]
Para todo $n\in\N_0$, o intervalo $I_n^x$ definido em \eqref{eq:18} satisfaz
\[
I_n^x=J_{k(n),m(n)}.
\]
\end{lema}
\begin{proof}
Como $b_n\ne0$, $\ell_n$ é uma bijeção afim estritamente monótona. Sua restrição a $[x_L,x_R]$ tem imagem $[c_n,d_n]$ por \eqref{eq:16}; logo $\ell_n^{-1}([c_n,d_n])=[x_L,x_R]$. Os dois quocientes em \eqref{eq:18} são precisamente $\ell_n^{-1}(c_n)$ e $\ell_n^{-1}(d_n)$, possivelmente em ordem invertida quando $b_n<0$.
\end{proof}
\begin{teorema}[unicidade do endereço]
$\Gamma$ é injetiva. Portanto, é uma bijeção de $\N_0$ sobre $\LL$ e cada átomo tem um único endereço.
\end{teorema}
\begin{proof}
Suponha $\Gamma(n)=\Gamma(n')=(a,b,[c,d])$. Pelo lema~3.4, a tripla recupera a célula comum $J=[x_L,x_R]$. Seu comprimento determina unicamente $k=-\log_2|J|$, e $m=2^kx_L$. Assim, também se recupera $x_0=(m+1/2)2^{-k}$. A inclinação local é $s=b$ e o valor ancorado é $y_0=a+bx_0$.

Pelas bijeções dos lemas~2.2, 2.4 e 3.1, os dados $(x_0,y_0,s)$ determinam unicamente $(k,\mu,r,t)$. O lema~3.2 determina então um único $n$. Logo $n=n'$.
\end{proof}
Em uma fronteira, duas células fechadas podem estar ativas. A célula semiaberta escolhe uma proprietária; em uma montagem contínua, as duas peças têm o mesmo valor no nó.
\begin{teorema}[dicionário de jatos racionais]
O mapa
\eqn{19}{n\longmapsto(x_{0,n},y_{0,n},s_n)}
é uma bijeção de $\N_0$ sobre $\D_{\rm mid}\times\Q\times\Q^*$. Sua imagem é densa em $\R^3$.
\end{teorema}
\begin{proof}
A bijetividade resulta da composição das bijeções em \eqref{eq:13}--\eqref{eq:14}. A densidade decorre da densidade de $\D_{\rm mid}$, de $\Q$ e de $\Q^*$ em $\R$, aplicada coordenada a coordenada na norma do máximo.
\end{proof}
O dicionário \eqref{eq:19} parametriza todos os 1-jatos com centro diádico não inteiro, valor racional e derivada racional não nula. Montagens contínuas ligam esses dados locais por valores de fronteira compartilhados.


% --- INICIO SECAO 4: ate Teorema 4.2 e equacao 23 (transcricao literal) ---
\section{Montagens e densidade em espaços de funções contínuas}
\begin{definicao}[montagem contínua]
Se $K=[\alpha,\beta]$ é um intervalo compacto não degenerado, $\Asm(K)$ é o conjunto das funções $g\in C(K)$ para as quais existe uma coleção finita $\Lambda$ de células diádicas, com interiores disjuntos e união contendo $K$, e endereços $n(J)$, $J\in\Lambda$, tais que
\eqn{20}{g|_{J\cap K}=\ell_{n(J)}|_{J\cap K},}
com valores coincidentes nas fronteiras compartilhadas.
\end{definicao}
Para $f\in C(K)$, denote o módulo de continuidade por
\eqn{21}{\omega(f;\delta)=\sup\{|f(x)-f(y)|:x,y\in K,\ |x-y|\le\delta\}.}
O resultado seguinte impõe simultaneamente nós diádicos, valores racionais e inclinações não nulas. A figura~\ref{fig:montagem} mostra a ligação das peças por nós comuns.
\figura{03_montagem_continua}{1}{Montagem contínua de átomos. À esquerda, cada cor identifica uma célula e seu segmento; peças consecutivas compartilham o mesmo valor no nó. À direita, o refinamento substitui células por filhos e aproxima a função-alvo com mais peças conectadas.}{montagem}

\begin{teorema}[densidade das montagens]
Para todo $f\in C[\alpha,\beta]$ e todo $\varepsilon>0$, existe $g\in\Asm([\alpha,\beta])$ tal que
\[
\normK{f-g}{[\alpha,\beta]}<\varepsilon.
\]
Todas as peças de $g$ pertencem a $\LL$ e têm inclinação não nula.
\end{teorema}
\begin{proof}
Estenda $f$ por constantes: $F(x)=f(\alpha)$ para $x<\alpha$, $F=f$ em $[\alpha,\beta]$ e $F(x)=f(\beta)$ para $x>\beta$. A extensão é uniformemente contínua. Escolha $k$ de modo que, com $h=2^{-k}$,
\eqn{22}{\omega(F;h)<\varepsilon/2.}
Tome os nós $t_j=jh$ desde $j_{\min}=\lfloor\alpha/h\rfloor$ até $j_{\max}=\lceil\beta/h\rceil$. Para cada $j$, escolha racional $v_j$ com $|v_j-F(t_j)|<\varepsilon/2$ e, sequencialmente, imponha $v_{j+1}\ne v_j$. Isso é possível porque cada intervalo aberto de escolha contém infinitos racionais e apenas um valor é proibido em cada passo.

Seja $g$ o interpolante linear dos pares $(t_j,v_j)$. Em $J_{k,j}$, sua inclinação é $(v_{j+1}-v_j)/h\in\Q^*$ e seu valor no centro é $(v_j+v_{j+1})/2\in\Q$. Pelo teorema~3.6, a peça é o único átomo com esses parâmetros. Os valores nodais compartilhados tornam $g$ contínua.

Se $P_kF$ é o interpolante dos valores exatos, então, para $x=(1-\theta)t_j+\theta t_{j+1}$,
\[
|F(x)-P_kF(x)|\le(1-\theta)|F(x)-F(t_j)|+\theta|F(x)-F(t_{j+1})|
\le\omega(F;h)<\varepsilon/2.
\]
Além disso, $P_kF-g$ é linear em cada célula e seus valores nodais têm módulo menor que $\varepsilon/2$, logo $\norm{P_kF-g}<\varepsilon/2$. A desigualdade triangular conclui.
\end{proof}
O número de células dessa cobertura uniforme é
\eqn{23}{N_k=\lceil\beta2^k\rceil-\lfloor\alpha2^k\rfloor=(\beta-\alpha)2^k+O(1).}


% --- INICIO DA SECAO 5 (texto original); preservar numeros (24)--(29) reservados na secao 4 omitida ---
\setcounter{equation}{29}
\section{Taxas uniformes e constante assintótica}
Para $f\in C^2[\alpha,\beta]$, use a extensão
\eqn{30}{F(x)=\begin{cases}
f(\alpha)+f'(\alpha)(x-\alpha)+\frac12f''(\alpha)(x-\alpha)^2,&x<\alpha,\\
f(x),&\alpha\le x\le\beta,\\
f(\beta)+f'(\beta)(x-\beta)+\frac12f''(\beta)(x-\beta)^2,&x>\beta.
\end{cases}}
Então $F\in C^2(\R)$ e $\normK{F''}{\R}=\normK{f''}{[\alpha,\beta]}$. Seja $P_kF$ o interpolante nos nós $j2^{-k}$.

\begin{teorema}[cota aguda e limite]
Ponha $h=2^{-k}$ e $M=\normK{f''}{[\alpha,\beta]}$. Então
\eqn{31}{\normK{f-P_kF}{[\alpha,\beta]}\le\frac{h^2M}8,}
e
\eqn{32}{\lim_{k\to\infty}\frac8{h^2}\normK{f-P_kF}{[\alpha,\beta]}=M.}
Além disso, para todo $\delta>0$ existe $g_{k,\delta}\in\Asm([\alpha,\beta])$ com
\eqn{33}{\norm{f-g_{k,\delta}}\le\frac{h^2M}8+\delta.}
\end{teorema}
\begin{proof}
Em uma célula $[x_L,x_R]$ de largura $h$, ponha $e=F-P_kF$. Para $x\in(x_L,x_R)$, a função
\[
\psi(t)=e(t)-\frac{(t-x_L)(t-x_R)}{(x-x_L)(x-x_R)}e(x)
\]
anula-se em $x_L,x,x_R$. Três aplicações de Rolle produzem $\xi_x\in(x_L,x_R)$ com $\psi''(\xi_x)=0$; portanto
\eqn{34}{F(x)-P_kF(x)=\frac12F''(\xi_x)(x-x_L)(x-x_R).}
O módulo do produto é no máximo $h^2/4$, com igualdade no centro, o que prova \eqref{eq:31}.

Se $M=0$, \eqref{eq:32} é imediata. Se $M>0$, escolha $x_*\in[\alpha,\beta]$ com $|f''(x_*)|=M$. Para cada $k$ suficientemente grande, escolha um centro de escala $k$, $z_k\in[\alpha,\beta]$, com $z_k\to x_*$. Avaliando \eqref{eq:34} em $z_k$,
\[
|F(z_k)-P_kF(z_k)|=\frac{h^2}8|F''(\xi_k)|,
\]
onde $\xi_k$ pertence à célula de $z_k$. O diâmetro da célula tende a zero, logo $\xi_k\to x_*$ e $|F''(\xi_k)|\to M$. Isso dá o limite inferior; \eqref{eq:31} fornece o superior.

Para a última afirmação, racionalize todos os valores nodais da cobertura finita com erro menor que $\delta$ e mantenha valores adjacentes distintos, como no teorema~4.2. O interpolante resultante é uma montagem e dista menos que $\delta$ de $P_kF$ em norma uniforme. A desigualdade triangular conclui.
\end{proof}

\begin{corolario}[preservação da constante]
Existe uma sequência $g_k\in\Asm([\alpha,\beta])$ tal que
\[
\lim_{k\to\infty}\frac8{h^2}\norm{f-g_k}=M.
\]
\end{corolario}
\begin{proof}
Na construção do teorema~5.1, escolha erros nodais $\delta_k=o(h^2)$. Então
\[
\left|\norm{f-g_k}-\norm{f-P_kF}\right|\le\norm{g_k-P_kF}<\delta_k.
\]
Divida por $h^2/8$ e aplique \eqref{eq:32}.
\end{proof}
Com $N_k$ de \eqref{eq:23}, $h=(\beta-\alpha)N_k^{-1}(1+O(N_k^{-1}))$ e, portanto,
\eqn{35}{\norm{f-P_kF}\le\frac{(\beta-\alpha)^2M}8N_k^{-2}\bigl(1+O(N_k^{-1})\bigr).}
Quando os extremos já são nós da malha, o termo de borda desaparece.
\begin{table}[!htbp]\centering
\caption{Erro exato do interpolante de $e^x$ em $[0,1]$.}
\begin{tabular}{rrrr}\toprule
$k$&erro exato&$2^{-2k}e/8$&razão\\\midrule
2&$1{,}8774\times10^{-2}$&$2{,}1237\times10^{-2}$&$0{,}8840$\\
4&$1{,}2866\times10^{-3}$&$1{,}3273\times10^{-3}$&$0{,}9693$\\
6&$8{,}2310\times10^{-5}$&$8{,}2955\times10^{-5}$&$0{,}9922$\\
8&$5{,}1746\times10^{-6}$&$5{,}1847\times10^{-6}$&$0{,}9980$\\
10&$3{,}2389\times10^{-7}$&$3{,}2404\times10^{-7}$&$0{,}9995$\\\bottomrule
\end{tabular}
\end{table}
\figura{05_taxa_exponencial}{.98}{Erro do interpolante de $e^x$ e razão entre erro e cota. A razão converge a $1$, conforme \eqref{eq:32}.}{taxa}


% --- INICIO DA SECAO 6 (texto original; figura externa substituida por legenda) ---
\section{Refinamento diádico adaptativo certificado}
O refinamento usa majorantes locais validados de $|F''|$ para controlar o erro em cada célula.

\subsection{Fechamento operacional por aritmética intervalar}
Seja $\EE$ a menor classe de expressões que contém $x$ e as constantes racionais e é fechada por soma, subtração, produto, recíproco, $\exp$, $\sin$, $\cos$ e composição. Um recíproco é admissível em intervalos nos quais seu denominador não se anula. Essa classe contém os polinômios, as funções racionais sem polos no domínio e as funções elementares usadas nos exemplos.

Composição significa substituição sintática: em $u\circ v$, cada folha $x$ da árvore de $u$ é substituída pela árvore de $v$. Toda expressão composta é, assim, uma árvore das operações-base.

Escreva $\mathbb I_{\Q}$ para os intervalos fechados com extremos racionais. A rotina
\eqn{36}{\operatorname{BoundSecondDerivative}(E,J;P,d)
=[\underline m_J,\overline m_J]\in\mathbb I_{\Q}}
usa precisão transcendental $P$ e profundidade de subdivisão $d$. Ela propaga simultaneamente intervalos para $(E,E',E'')$, isto é, 2-jatos intervalares. Por exemplo,
\begin{align*}
(uv)''&=u''v+2u'v'+uv'',\\
(u^{-1})''&=2(u')^2u^{-3}-u''u^{-2},\\
(e^u)''&=e^u\bigl((u')^2+u''\bigr),\\
(\sin u)''&=-\sin(u)(u')^2+\cos(u)u'',\\
(\cos u)''&=-\cos(u)(u')^2-\sin(u)u''.
\end{align*}
As operações racionais são exatas. Para $\exp$, a redução do argumento a $|z|\le1/2$ é seguida por uma soma de Taylor com resto majorado por uma série geométrica racional; para $\sin$ e $\cos$, usa-se o resto de Lagrange racional e a interseção com $[-1,1]$. Elevar $P$ e $d$ reduz a sobrestimação.

\begin{lema}[restos transcendentais com extremos racionais]
Se $I=[a,b]\in\mathbb I_{\Q}$, as rotinas descritas acima produzem intervalos de extremos racionais que contêm $\exp(I)$, $\sin(I)$ e $\cos(I)$. Mais explicitamente, para $0\le z\le1/2$ e $n\ge0$,
\eqn{37}{0\le e^z-\sum_{j=0}^n\frac{z^j}{j!}
\le\frac{z^{n+1}}{(n+1)!}\frac1{1-z/(n+2)}.}
Se $R=\max\{|a|,|b|\}$ e
\[
P_n^{\sin}(z)=\sum_{j=0}^n(-1)^j\frac{z^{2j+1}}{(2j+1)!},\qquad
P_n^{\cos}(z)=\sum_{j=0}^n(-1)^j\frac{z^{2j}}{(2j)!},
\]
então, para todo $z\in I$,
\eqn{38}{|\sin z-P_n^{\sin}(z)|\le\frac{R^{2n+2}}{(2n+2)!},\qquad
|\cos z-P_n^{\cos}(z)|\le\frac{R^{2n+1}}{(2n+1)!}.}
Avaliar os polinômios em $I$, inflar pelos respectivos restos e intersectar com $[-1,1]$ preserva a inclusão.
\end{lema}
\begin{proof}
Na série de $e^z$, o primeiro termo omitido é $t_{n+1}=z^{n+1}/(n+1)!$. Para $j\ge n+1$, a razão entre termos consecutivos é $z/(j+1)\le z/(n+2)<1$; a cauda é, portanto, majorada pela série geométrica de primeiro termo $t_{n+1}$ e razão $z/(n+2)$, o que dá \eqref{eq:37}. Para um argumento racional $x$ arbitrário, escolha $r$ com $z=|x|/2^r\le1/2$, aplique a inclusão anterior e eleve o intervalo ao quadrado $r$ vezes; se $x<0$, tome o recíproco do intervalo positivo obtido para $-x$. Em $I=[a,b]$, a monotonicidade de $e^x$ permite usar o extremo inferior da inclusão em $a$ e o superior da inclusão em $b$.

As duas desigualdades em \eqref{eq:38} são o resto de Lagrange dos polinômios de Maclaurin de graus $2n+1$ e $2n$, pois todas as derivadas de seno e cosseno têm módulo no máximo $1$. A avaliação intervalar contém a imagem de cada polinômio; inflar pelo raio do resto contém a função. Como seno e cosseno já pertencem a $[-1,1]$, a interseção não remove nenhum valor possível. Todas as operações, fatoriais e extremos envolvidos são racionais.
\end{proof}

\begin{proposicao}[correção do majorante intervalar]
Se $E\in\EE$ define uma função $C^2$ em $J$ e a rotina \eqref{eq:36} termina com sucesso, então
\eqn{39}{E''(x)\in[\underline m_J,\overline m_J]\ (x\in J),\qquad
M_{\rm int}(J):=\max\{|\underline m_J|,|\overline m_J|\}\ge\normK{E''}{J}.}
Se algum intervalo de denominador contém zero mesmo após a subdivisão permitida, a rotina devolve falha e nenhum número é aceito como certificado.
\end{proposicao}
\begin{proof}
Para constantes e para $x$, a inclusão é imediata. Soma, produto e recíproco preservam inclusão pelas operações elementares de intervalos; o recíproco só é executado quando zero está fora do intervalo do denominador. As fórmulas de 2-jatos acima e o lema~6.1 dão o passo de indução estrutural para toda expressão de $\EE$. Pela normalização sintática declarada antes do lema, uma composição genérica já foi expandida nessa árvore e não acrescenta um caso independente. Subdividir $J$ preserva a cobertura, e tomar o fecho convexo dos intervalos obtidos em suas partes contém a união das imagens de $E''$. Isso prova a primeira inclusão e, tomando módulos, \eqref{eq:39}. A regra de falha impede que uma divisão não validada produza uma falsa garantia.
\end{proof}
A propagação registra a largura dos intervalos, o número de caixas e a profundidade usados. Subdivisão e Taylor racional permitem apertar as inclusões.

\begin{algoritmo}[refinamento certificado]
Dados $f\in C^2[\alpha,\beta]$, tolerâncias $\varepsilon_I,\varepsilon_R>0$ e a extensão $F$ de \eqref{eq:30}, use, para toda célula $J$, um majorante validado
\eqn{40}{\normK{F''}{J}\le M(J)\le M_*<\infty.}
Quando $F$ é fornecida por uma expressão de $\EE$, tome $M(J)=M_{\rm int}(J)$ de \eqref{eq:39}; se a certificação falhar, aborte sem emitir uma montagem garantida. Para uma função $C^2$ fora dessa interface, \eqref{eq:40} é um contrato de entrada.
\begin{enumerate}[label=(\roman*),topsep=6pt,itemsep=3pt]
\item Inicie com uma cobertura diádica finita de $[\alpha,\beta]$.
\item Para cada folha $J$, calcule
\eqn{41}{\eta_{\rm cert}(J)=\frac{|J|^2M(J)}8.}
\item Enquanto existir folha com $\eta_{\rm cert}(J)>\varepsilon_I$, substitua-a por seus dois filhos.
\item Nos nós finais, escolha racionais compartilhados $v_t$ com $|v_t-F(t)|<\varepsilon_R$ e valores adjacentes distintos.
\item Emita, em cada folha, o endereço do átomo que interpola seus dois valores nodais.
\end{enumerate}
\end{algoritmo}

\begin{teorema}[correção, término e tamanho]
O algoritmo~6.3 termina e produz $g\in\Asm([\alpha,\beta])$ com
\eqn{42}{\norm{f-g}\le\varepsilon_I+\varepsilon_R.}
Se a cobertura inicial tem $R_0$ raízes de largura máxima $h_0$ e $M_*>0$, nenhuma ramificação ultrapassa
\eqn{43}{r_*=\max\left\{0,\left\lceil\frac12\log_2\left(\frac{h_0^2M_*}{8\varepsilon_I}\right)\right\rceil\right\}}
níveis, e há no máximo $R_02^{r_*}$ folhas. Se $M_*=0$, pode-se tomar $r_*=0$.
\end{teorema}
\begin{proof}
Após $r$ subdivisões, uma folha descendente de uma raiz tem $|J|\le h_02^{-r}$ e, por \eqref{eq:40},
\[
\eta_{\rm cert}(J)\le\frac{h_0^2M_*}8\,4^{-r}.
\]
O lado direito é no máximo $\varepsilon_I$ para $r\ge r_*$; logo nenhuma folha é subdividida além dessa profundidade. A árvore é finita e cada raiz gera no máximo $2^{r_*}$ folhas.

Na partição final, \eqref{eq:34} e \eqref{eq:40} implicam
\[
\normK{F-P_\Lambda F}{J}\le\frac{|J|^2}8\normK{F''}{J}\le\eta_{\rm cert}(J)\le\varepsilon_I.
\]
A racionalização nodal produz, pelos teoremas~3.5 e 3.6, uma montagem de átomos de endereço único e com $\norm{P_\Lambda F-g}<\varepsilon_R$. Como $F=f$ no intervalo original, a desigualdade triangular prova \eqref{eq:42}.
\end{proof}
\figura{06_refinamento_certificado}{.90}{Refinamento certificado para $f(x)=[1+64(x-5/8)^2]^{-1}$. Células com certificado acima de $\varepsilon_I$ são subdivididas. A malha se concentra na região de maior curvatura, e o painel inferior compara os certificados finais com a tolerância. A faixa superior representa o limite uniforme $\varepsilon_I+\varepsilon_R$.}{refinamento}


\begin{proposicao}[calibração exata para curvatura constante]
Suponha que $[\alpha,\beta]=J_0$ seja uma única célula diádica inicial, $F''\equiv\kappa\ne0$ em $J_0$ e $M(J)=|\kappa|$ em todo descendente. Pondo
\[
A=\frac{|J_0|^2|\kappa|}8,
\]
o número de folhas produzido com tolerância $\varepsilon_I$ é
\eqn{44}{N(\varepsilon_I)=2^{r(\varepsilon_I)},\qquad
r(\varepsilon_I)=\max\left\{0,\left\lceil\log_4\left(\frac A{\varepsilon_I}\right)\right\rceil\right\}.}
Consequentemente,
\eqn{45}{\begin{aligned}
\liminf_{\varepsilon_I\downarrow0}\sqrt{\varepsilon_I}\,N(\varepsilon_I)
&=\frac1{\sqrt8}\int_\alpha^\beta|f''(x)|^{1/2}\,dx,\\
\limsup_{\varepsilon_I\downarrow0}\sqrt{\varepsilon_I}\,N(\varepsilon_I)
&=\frac1{\sqrt2}\int_\alpha^\beta|f''(x)|^{1/2}\,dx.
\end{aligned}}
A quantização diádica introduz, nesse caso, um fator dois entre os limites de acumulação.
\end{proposicao}
\begin{proof}
Após $r$ subdivisões uniformes, cada folha tem largura $|J_0|2^{-r}$ e certificado $A4^{-r}$. Como todos os certificados de um mesmo nível são iguais, o algoritmo para no primeiro $r$ com $A4^{-r}\le\varepsilon_I$, o que prova \eqref{eq:44}. Para $A4^{-r}\le\varepsilon_I<A4^{-(r-1)}$, tem-se
\[
\sqrt A\le\sqrt{\varepsilon_I}\,2^r<2\sqrt A.
\]
Os extremos inferior e superior são atingidos, respectivamente, nos limiares $\varepsilon_I=A4^{-r}$ e por aproximação pela esquerda aos limiares $A4^{-(r-1)}$. Por fim, $\int_\alpha^\beta|f''|^{1/2}=|J_0|\sqrt{|\kappa|}$.
\end{proof}

\begin{observacao}[diferença segunda e erro no centro]
Se $x_0=(x_L+x_R)/2$, então
\[
|f(x_0)-P_Jf(x_0)|=\frac12|f(x_L)-2f(x_0)+f(x_R)|.
\]
Isso mede o erro no centro. Para $f(x)=x^3$ em $[-1,1]$, $P_Jf(x)=x$ e a diferença segunda central é zero, enquanto
\[
\norm{x^3-x}=\frac2{3\sqrt3}\approx0{,}3849.
\]
O controle em todo o intervalo é fornecido pelo majorante de \eqref{eq:40}.
\end{observacao}

\begin{thebibliography}{1}
\bibitem[Hardy e Wright(2008)]{hardy2008} Hardy, G. H.; Wright, E. M. \emph{An Introduction to the Theory of Numbers}, 6a ed., Oxford University Press, 2008.
\end{thebibliography}
\end{document}
